\section{Resistències en paral·lel}

Dues resistències estan \textbf{en paral·lel} quan els dos extrems de cadascuna estan connectats als mateixos dos punts: el corrent té \emph{dos camins} per escollir. Pensa en un tub d'aigua que es divideix en dos tubs i després torna a ajuntar-se: a més camins, més aigua pot passar.

\begin{center}\cparalel{$V_T$}{$R_1$}{$R_2$}\end{center}

\begin{idea}
En un circuit en paral·lel:
\begin{enumerate}[label=\textbf{\arabic*.}]
 \item La \textbf{tensió és la mateixa} a totes les branques: \quad $V_T = V_1 = V_2 = \cdots$
 \item La \textbf{intensitat es reparteix}: \quad $I_T = I_1 + I_2 + \cdots$
 \item La \textbf{resistència total} es calcula així: \quad $\dfrac{1}{R_T} = \dfrac{1}{R_1} + \dfrac{1}{R_2} + \cdots$
\end{enumerate}
\end{idea}

\begin{idea}
\textbf{Dos casos que s'usen molt:}
\begin{itemize}
 \item \textbf{Dues resistències:} $R_T = \dfrac{R_1 \cdot R_2}{R_1 + R_2}$
 \item \textbf{$n$ resistències iguals de valor $R$:} $R_T = \dfrac{R}{n}$
\end{itemize}
\textbf{Regla per comprovar:} la resistència total en paral·lel \textbf{sempre és més petita} que la més petita de les resistències.
\end{idea}

\begin{idea}
\textbf{Càlcul ràpid amb calculadora científica (tecla $x^{-1}$ o $1/x$):}
\\[0.3em]
Pots calcular la resistència total directament aïllant la fórmula com:
\[
  R_T = \left( \frac{1}{R_1} + \frac{1}{R_2} + \cdots \right)^{-1}
\]
\textbf{Passos a seguir:}
\begin{enumerate}[label=\textbf{\arabic*.}]
 \item Suma les inverses de cada resistència: \quad $R_1\ \mathbf{x^{-1}} + R_2\ \mathbf{x^{-1}} + \dots =$
 \item Inverteix el resultat obtingut: \quad $\mathbf{x^{-1}} =$ \quad (o bé $\text{Ans}^{-1} =$)
\end{enumerate}
\end{idea}

\begin{metode}{calcular un circuit en paral·lel}
\begin{itemize}
 \item\textbf{\ding{172}} Calcula $R_T$ amb la fórmula. \quad
 \item\textbf{\ding{173}} Calcula la intensitat total: $I_T = V_T / R_T$. \quad
 \item\textbf{\ding{174}} Calcula la intensitat de cada branca (la tensió és la mateixa): $I_i = V_T / R_i$. \quad
 \item\textbf{\ding{175}} Calcula les potències: $P_i = V_T \cdot I_i$. \quad
 \item\textbf{\ding{176}} \textbf{Comprova}: les intensitats sumen $I_T$? les potències sumen $V_T \cdot I_T$?
\end{itemize}
\end{metode}

\begin{exemple}{dues resistències en paral·lel}
Una pila de \SI{12}{V} alimenta $R_1 = \SI{200}{\ohm}$ i $R_2 = \SI{300}{\ohm}$ en paral·lel.
\begin{center}\cparalel{\SI{12}{V}}{$R_1=\SI{200}{\ohm}$}{$R_2=\SI{300}{\ohm}$}\end{center}
\begin{itemize}
 \item $R_T = \dfrac{200 \cdot 300}{200 + 300} = \dfrac{60000}{500} = \SI{120}{\ohm}$ \ (menor que \SI{200}{\ohm}: té sentit)
 \item $I_T = 12 / 120 = \SI{0,1}{A}$
 \item $I_1 = 12 / 200 = \SI{0,06}{A}$ \qquad $I_2 = 12 / 300 = \SI{0,04}{A}$
 \item $P_1 = 12 \cdot 0{,}06 = \SI{0,72}{W}$ \qquad $P_2 = 12 \cdot 0{,}04 = \SI{0,48}{W}$ \qquad $P_T = 12 \cdot 0{,}1 = \SI{1,2}{W}$
 \item \textbf{Comprovació:} $0{,}06 + 0{,}04 = 0{,}1$ \checkmark \qquad $0{,}72 + 0{,}48 = 1{,}2$ \checkmark
\end{itemize}
\begin{taulamag}
$R_1$ & 12 & 0{,}06 & 200 & 0{,}72\\
$R_2$ & 12 & 0{,}04 & 300 & 0{,}48\\
\midrule
Total & 12 & 0{,}1 & 120 & 1{,}2\\
\end{taulamag}
\end{exemple}

\begin{exemple}{resistències iguals}
Tres resistències de \SI{90}{\ohm} en paral·lel: $R_T = \dfrac{90}{3} = \SI{30}{\ohm}$.
\end{exemple}

\begin{atencio}
\begin{itemize}
 \item \textbf{(1)} Sumar les resistències directament, com en sèrie. En paral·lel \emph{no} es sumen. \quad
 \item \textbf{(2)} Oblidar l'últim pas: si has obtingut $\dfrac{1}{R_T} = 0{,}05$, encara no has acabat: $R_T = \dfrac{1}{0{,}05} = \SI{20}{\ohm}$.
\end{itemize}
\end{atencio}

\clearpage
\begin{guiada}{pila de \SI{6}{V} amb $R_1 = \SI{100}{\ohm}$ i $R_2 = \SI{300}{\ohm}$ en paral·lel}
\pas{1}{Resistència total:}{\vspace{0.7cm}}
\pas{2}{Intensitat total:}{\vspace{0.7cm}}
\pas{3}{Intensitat de cada branca:}{\vspace{0.7cm}}
\begin{taulamag}
$R_1$ & 6 & & 100 &\\
$R_2$ & 6 & & 300 &\\
\midrule
Total & 6 & & &\\
\end{taulamag}
\resposta{0pt}{$R_T = \dfrac{100 \cdot 300}{400} = \SI{75}{\ohm}$; $I_T = 6/75 = \SI{0,08}{A}$; $I_1 = 6/100 = \SI{0,06}{A}$; $I_2 = 6/300 = \SI{0,02}{A}$.
\begin{taulamag}
$R_1$ & 6 & 0{,}06 & 100 & 0{,}36\\
$R_2$ & 6 & 0{,}02 & 300 & 0{,}12\\
\midrule
Total & 6 & 0{,}08 & 75 & 0{,}48\\
\end{taulamag}}
\end{guiada}
\clearpage
\subsection*{Exercicis}

\begin{exercici}{1}
Calcula la resistència total en paral·lel de:\quad
\textbf{a)} dues resistències de \SI{100}{\ohm}.\qquad
\textbf{b)} tres resistències de \SI{60}{\ohm}.\qquad
\textbf{c)} \SI{40}{\ohm} i \SI{60}{\ohm}.
\resposta{2.4cm}{a) $100/2 = \SI{50}{\ohm}$. \quad b) $60/3 = \SI{20}{\ohm}$. \quad c) $\dfrac{40 \cdot 60}{100} = \SI{24}{\ohm}$.}
\end{exercici}

\begin{exercici}{1}
Calcula la resistència total en paral·lel de:\quad
\textbf{a)} \SI{12}{\ohm} i \SI{4}{\ohm}.\qquad
\textbf{b)} \SI{10}{\ohm} i \SI{40}{\ohm}.
\resposta{2.4cm}{a) $\dfrac{12 \cdot 4}{16} = \SI{3}{\ohm}$. \quad b) $\dfrac{10 \cdot 40}{50} = \SI{8}{\ohm}$.}
\end{exercici}

\begin{exercici}[20]{2}
Completa la taula del circuit següent.
\begin{center}\cparalel{\SI{12}{V}}{$R_1=\SI{40}{\ohm}$}{$R_2=\SI{60}{\ohm}$}\end{center}
\begin{taulamag}
$R_1$ & 12 & & 40 &\\
$R_2$ & 12 & & 60 &\\
\midrule
Total & 12 & & &\\
\end{taulamag}
\resposta{0pt}{$R_T = \SI{24}{\ohm}$, $I_T = \SI{0,5}{A}$.
\begin{taulamag}
$R_1$ & 12 & 0{,}3 & 40 & 3{,}6\\
$R_2$ & 12 & 0{,}2 & 60 & 2{,}4\\
\midrule
Total & 12 & 0{,}5 & 24 & 6\\
\end{taulamag}}
\end{exercici}

\begin{exercici}[20]{2}
Completa la taula del circuit següent, amb tres resistències en paral·lel.\quad (Pista: $\frac{1}{R_T} = \frac{1}{100} + \frac{1}{200} + \frac{1}{200}$.)
\begin{center}\cparaleltres{\SI{10}{V}}{$R_1=\SI{100}{\ohm}$}{$R_2=\SI{200}{\ohm}$}{$R_3=\SI{200}{\ohm}$}\end{center}
\begin{taulamag}
$R_1$ & 10 & & 100 &\\
$R_2$ & 10 & & 200 &\\
$R_3$ & 10 & & 200 &\\
\midrule
Total & 10 & & &\\
\end{taulamag}
\resposta{0pt}{$\dfrac{1}{R_T} = 0{,}01 + 0{,}005 + 0{,}005 = 0{,}02$, per tant $R_T = \SI{50}{\ohm}$; $I_T = \SI{0,2}{A}$.
\begin{taulamag}
$R_1$ & 10 & 0{,}1 & 100 & 1\\
$R_2$ & 10 & 0{,}05 & 200 & 0{,}5\\
$R_3$ & 10 & 0{,}05 & 200 & 0{,}5\\
\midrule
Total & 10 & 0{,}2 & 50 & 2\\
\end{taulamag}}
\end{exercici}

\begin{exercici}{3}
Una resistència de \SI{60}{\ohm} es connecta en paral·lel amb una altra resistència $R_2$, i la resistència total és de \SI{40}{\ohm}. Quant val $R_2$?
\resposta{3cm}{$\dfrac{1}{R_2} = \dfrac{1}{40} - \dfrac{1}{60} = \dfrac{3 - 2}{120} = \dfrac{1}{120}$, per tant $R_2 = \SI{120}{\ohm}$. Comprovació: $\dfrac{60 \cdot 120}{180} = \SI{40}{\ohm}$ \checkmark}
\end{exercici}

\begin{exercici}{3}
Dues resistències $R_1$ i $R_2$ en paral·lel estan connectades a \SI{6}{V}. Un amperímetre indica que la intensitat total és \SI{0,5}{A} i que per $R_1$ hi passa \SI{0,2}{A}. Calcula $I_2$, $R_1$, $R_2$ i $R_T$.
\resposta{3.5cm}{$I_2 = 0{,}5 - 0{,}2 = \SI{0,3}{A}$. $R_1 = 6/0{,}2 = \SI{30}{\ohm}$. $R_2 = 6/0{,}3 = \SI{20}{\ohm}$. $R_T = 6/0{,}5 = \SI{12}{\ohm}$. Comprovació: $\dfrac{30 \cdot 20}{50} = \SI{12}{\ohm}$ \checkmark}
\end{exercici}

\begin{comprova}
A casa, els endolls estan connectats en paral·lel. Per què? Què passaria si estiguessin en sèrie?
\resposta{2.4cm}{En paral·lel, cada aparell rep la tensió completa (\SI{230}{V}) i funciona independentment dels altres. En sèrie, la tensió es repartiria i, si un aparell s'apagués, s'apagarien tots.}
\end{comprova}
